\documentclass[11pt]{article}
\usepackage[margin=0.82in]{geometry}
\usepackage{amsmath,amssymb,booktabs,array,longtable}
\usepackage[T1]{fontenc}
\usepackage{lmodern}
\usepackage{microtype}
\usepackage{xcolor}
\usepackage{hyperref}
\usepackage{enumitem}
\hypersetup{colorlinks=true,linkcolor=blue!55!black,urlcolor=blue!55!black,citecolor=blue!55!black}

\title{\textbf{One Mathematics of Arrays from Transformer Semantics to Machine Execution}\\
\large Shape, Indexing, Normal Forms, Machine Translation, Prediction, and Cross-Validation}
\author{Lenore M. Mullin\\
College of Nanotechnology, Science, and Engineering\\
University at Albany, SUNY, Albany, NY\\
\texttt{lmullin@albany.edu}}
\date{Working derivation -- October 2026}

\newcommand{\DNF}{\mathrm{DNF}}
\newcommand{\ONF}{\mathrm{ONF}}
\newcommand{\rav}{\operatorname{rav}}
\newcommand{\SiLU}{\operatorname{SiLU}}
\newcommand{\Norm}{\operatorname{Norm}}
\newcommand{\Attn}{\operatorname{Attn}}
\newcommand{\FFN}{\operatorname{FFN}}
\newcommand{\redplus}{+_{\mathrm{red}}}

\begin{document}
\maketitle

\begin{abstract}
This working document unifies the current Mathematics of Arrays (MoA) Transformer derivation
with the larger experimental and compiler program.  The governing claim is deliberately simple:
there is one array mathematics from semantic data to shaped machines.  Arrays are defined by
shape; valid index spaces are generated from shape; components are selected by $\psi$; and
computation is reduced to a machine-independent Denotational Normal Form (DNF).  Only after the
DNF is fixed is a realization introduced through reshape/dimension lifting, a chosen member of
the $\gamma$ family, and \(\rav\).  Row-major $\gamma$ is used as the default concrete example,
but column-major and other layout functions are equally definable.  A change from one $\gamma$
to another is an explicit translation and its cost must be derived rather than hidden.

The same theory is required to explain and predict the independent CPU, GPU, multi-GPU,
distributed, and heterogeneous CPU--GPU experiments.  Those experiments are not separate
performance models.  They are independent tests of the same shape/index/mapping theory and
therefore must mutually corroborate its semantic, resource, communication, and predictive
consequences.  The Transformer is developed under the same discipline.  The existing MoA
attention-family derivations are treated as established components and are composed with
normalization, residual, and feed-forward definitions to obtain a complete block DNF before any
machine is introduced.
\end{abstract}

\section{One Theory, All the Way}

The governing research chain is
\[
\boxed{\begin{aligned}
\text{FORMULATE}&\rightarrow\text{DERIVE/OPTIMIZE}\rightarrow\DNF\\
&\xrightarrow{\gamma}\text{SEQUENTIAL ONF}\rightarrow\text{LIFTED ONF}\\
&\rightarrow\text{MACHINE ARRAY}\rightarrow\text{PREDICT/FREEZE}\rightarrow\text{EXECUTE/VALIDATE}.
\end{aligned}}
\]

The central invariant is
\[
\boxed{\textbf{The DNF never changes when the target machine changes.}}
\]

The algorithm, its data, its dependency structure, and the target machine are all viewed as
arrays.  The same MoA rules apply on both sides of the conventional hardware/software boundary.
The backend may be CPU, GPU, multi-GPU, MPI+OpenMP, OpenACC, CUDA, or concurrent CPU--GPU;
the mathematical semantics are not replaced by a backend-specific theory.

\subsection{The minimal semantic vocabulary}

For an array \(A\), its shape is
\[
\boxed{\rho A.}
\]
Its valid multidimensional index domain is generated from that shape:
\[
\boxed{\iota(\rho A).}
\]
For every valid index \(\vec i\in\iota(\rho A)\), the corresponding component is
\[
\boxed{\vec i\psi A.}
\]
The complete valid index structure reconstructs the array:
\[
\boxed{(\iota(\rho A))\psi A\equiv A.}
\]

The canonical catenation operator is \(\#\).  Dimensionality is scalar:
\[
\boxed{\delta(A)=\tau(\rho A).}
\]

These definitions do not depend on whether \(A\) represents data, a dependency object, a
partition, a processor array, a cache hierarchy, a set of GPUs, a collection of nodes, or another
machine property.

\subsection{The machine is an array too}

Let the target machine be \(M\), with
\[
\boxed{
\rho M=\langle m_0\;m_1\;\cdots\;m_{r-1}\rangle .
}
\]
Then
\[
\iota(\rho M)
\]
is its shaped machine-index space, and a valid machine component is selected by the same
operation
\[
\vec p\psi M,\qquad \vec p\in\iota(\rho M).
\]

Thus the common foundation is
\[
\boxed{
\begin{array}{c}
\text{DATA ARRAY}\\
\rho,\iota,\psi
\end{array}
\quad\longleftrightarrow\quad
\begin{array}{c}
\text{COMPUTATION ARRAY}\\
\rho,\iota,\psi
\end{array}
\quad\longleftrightarrow\quad
\begin{array}{c}
\text{MACHINE ARRAY}\\
\rho,\iota,\psi .
\end{array}}
\]

\section{DNF, Sequential ONF, and the \(\gamma\) Family}

The DNF contains the architecture-independent index/value computation.  No concrete machine
address belongs in the DNF.  The first operational realization is always the one-to-one sequential
ONF:
\[
\boxed{\DNF}
\xrightarrow[\text{one-to-one}]{\gamma}
\boxed{\text{Sequential ONF}}.
\]

\subsection{First use of \(\gamma\): storage-order convention}

The function \(\gamma\) translates a multidimensional index, together with an array shape, to a
linear address in a raveled representation.  Unless otherwise stated, concrete examples in this
document use a \textbf{row-major \(\gamma\)}.  This is a realization choice, not a restriction of
MoA.  A column-major \(\gamma\) is also defined, and other \(\gamma\)'s may be designed for tiled,
blocked, strided, sparse, hierarchical, or architecture-specific layouts.

For
\[
\rho A=\langle n\;m\rangle,
\]
the row-major member is
\[
\boxed{
\gamma_{\rm rm}(\langle i\#j\rangle;\rho A)
=
(i\times m)+j,
}
\]
and
\[
\boxed{
\langle i\#j\rangle\psi A
\longrightarrow
(\rav A)[(i\times m)+j].
}
\]

For the same abstract array, a column-major member is
\[
\boxed{
\gamma_{\rm cm}(\langle i\#j\rangle;\rho A)
=
i+(j\times n).
}
\]

Hence the DNF may have multiple legal operational realizations:
\[
\boxed{\DNF(F)}
\xrightarrow{\gamma_{\rm rm}}
\ONF_{\rm rm}(F),
\qquad
\boxed{\DNF(F)}
\xrightarrow{\gamma_{\rm cm}}
\ONF_{\rm cm}(F).
\]

\subsection{Changing \(\gamma\) has a cost}

If one stage is realized under \(\gamma_a\) and a later stage requires \(\gamma_b\), the change is
not silently free.  Either the new access can be interpreted without physical reorganization, or
an explicit translation is required:
\[
\boxed{
\gamma_a
\longrightarrow
T_{a\rightarrow b}
\longrightarrow
\gamma_b .
}
\]

The translation has a deterministic resource consequence
\[
\boxed{C(T_{a\rightarrow b})}
\]
that must be included in the mapped cost.  Thus layout becomes a mathematical design variable.
A predictive realization may compare
\[
\gamma_1,\gamma_2,\ldots,\gamma_k
\quad\longrightarrow\quad
C_1,C_2,\ldots,C_k
\]
before execution.

\section{Reshape and Dimension Lifting}

Dimension lifting occurs \emph{after} the sequential ONF has been formed and verified.  It does
not alter the DNF.

If a sequential execution index space has shape \(\rho P\), a legal lifted shape
\(\rho P^L\) must preserve the number of represented components:
\[
\boxed{\tau(\rho P^L)=\tau(\rho P).}
\]

A canonical lifting pattern is
\[
\boxed{N\rightarrow \Pi\#(N\div\Pi),}
\]
with the lifted application
\[
\boxed{
(f\Omega^{\langle1\rangle})
\left(
(\Pi\#(N\div\Pi))\rho N
\right).
}
\]

The semantic invariant is
\[
\boxed{
\vec i\psi Output_{\DNF}
=
\vec i\psi Output_{\rm seq}
=
\vec i\psi Output_{\rm lifted}.
}
\]

Therefore
\[
\boxed{\textbf{Dimension lifting changes execution shape, not semantic meaning.}}
\]

The lifted shape also supplies the legal loop/index bounds.  For
\[
\rho P^L=\langle p_0\;p_1\;\cdots\;p_{k-1}\rangle,
\]
the corresponding bounds are
\[
0\le i_0<p_0,\ldots,0\le i_{k-1}<p_{k-1}.
\]

\section{Prediction and Independent Experimental Cross-Validation}

The experimental discipline is
\[
\boxed{
\text{derive}
\rightarrow
\text{map}
\rightarrow
\text{predict}
\rightarrow
\textbf{FREEZE}
\rightarrow
\text{first execution}
\rightarrow
\text{validate}.
}
\]

The independent experiments must not become unrelated empirical models.  Each experiment
instantiates the same MoA theory with a different computation shape and/or machine shape.
Consequently, the experiments cross-validate the theory by testing different consequences of the
same definitions.

The separation is:
\begin{center}
\begin{tabular}{p{.28\textwidth}p{.62\textwidth}}
\textbf{Exact semantic quantities} & shapes, valid indices, dependencies, work, ownership, bytes, communication\\
\textbf{Machine realization properties} & capacities, rates, paths, latency, bandwidth, realization factors\\
\textbf{Validation quantities} & correctness, first execution, measured traffic/time, prediction error
\end{tabular}
\end{center}

No target measurement used for validation may be retroactively inserted into the frozen
prediction.

\subsection{What the CPU experiments teach}

The multicore CPU experiments instantiate shaped processors, cache domains, NUMA domains,
and memory resources.  The same computation DNF produces different realized behavior when its
ONF is mapped to different machine shapes.  Controlled compact-versus-balanced NUMA experiments
show that worker count alone is insufficient: the actual indexed machine shape and first-touch
placement matter.

At \(16384^2\), representative 1-to-64-core measurements include:
\[
\begin{array}{lrrrr}
\toprule
\text{Program} & T_1\;(\mathrm{s}) & T_{64}\;(\mathrm{s}) & \text{Speedup} & \text{Efficiency}\\
\midrule
\text{Sobel} & 1.178684 & 0.018703 & 63.02 & 98.5\%\\
\text{Reduce} & 0.366707 & 0.005836 & 62.84 & 98.2\%\\
\text{Histogram} & 0.157111 & 0.002637 & 59.59 & 93.1\%\\
\text{Pipeline} & 0.844528 & 0.015650 & 53.96 & 84.3\%\\
\text{Otsu} & 0.186954 & 0.003935 & 47.51 & 74.2\%\\
\text{RGB2Gray} & 0.105370 & 0.002781 & 37.89 & 59.2\%\\
\text{Threshold} & 0.030802 & 0.001728 & 17.82 & 27.8\%\\
\bottomrule
\end{array}
\]
The different scaling curves are consequences to be explained from different \(\psi\)-derived
dependency structures under the same machine.

\subsection{What the GPU and multi-GPU experiments teach}

GPU experiments retain the semantic DNF while replacing the machine array and operational
mapping.  Pointwise work, neighborhood work, reductions, shape-changing remaps, and composed
pipelines therefore become independent tests of the same theory.

For a row-partitioned pointwise array,
\[
\rho A=\langle H\;W\rangle,\qquad
\rho A_p=\left\langle H/P\;W\right\rangle,
\]
aligned ownership gives no semantic halo:
\[
\boxed{H_p=\varnothing.}
\]

For Sobel, the \(3\times3\) dependency causes required indices to cross an ownership boundary.
For a two-way row split the exact overlap is two rows:
\[
\boxed{N_{\rm halo}=2W.}
\]
The distinction is derived from indexing, not from a GPU-specific communication rule.

\subsection{What heterogeneous CPU--GPU experiments teach}

Concurrent CPU--GPU execution tests a stronger consequence: the same semantic index space may
be divided across unlike machine components.  The mapped completion time must account for the
distinct machine realization functions and any communication implied by the dependency structure.

For the pointwise Threshold experiment with \(N=2^{30}\), the frozen heterogeneous split was
\[
\boxed{
N_C=689{,}968{,}322,\qquad
N_G=383{,}773{,}502,
}
\]
with
\[
\boxed{\widehat T_{\rm Threshold}=86.2055821\ {\rm ms}.}
\]
The first measured completion was
\[
\boxed{T_{\rm measured}=85.816671\ {\rm ms},}
\]
giving
\[
\boxed{\text{prospective error}=-0.451\%.}
\]

A separately frozen heterogeneous Halide Sobel mapping predicted
\[
\boxed{\widehat T_{\rm Sobel}=6.495540\ {\rm ms}}
\]
and measured
\[
\boxed{T_{\rm measured}=6.465319\ {\rm ms},}
\]
for
\[
\boxed{\text{prospective error}=-0.4653\%.}
\]
The Sobel case additionally exposed a critical MoA point: nominal CPU identifiers did not initially
represent the intended physical-core machine shape.  Correcting the \emph{machine mapping}, while
leaving the frozen semantic design unchanged, restored correspondence between prediction and
execution.

These experiments therefore validate different parts of one chain:
\[
\boxed{
\text{indexing}
\rightarrow
\text{dependency}
\rightarrow
\text{ownership}
\rightarrow
\text{communication}
\rightarrow
\text{machine mapping}
\rightarrow
\text{prediction}.
}
\]

\subsection{Composition validates a resource tradeoff, not automatic speedup}

The composed RGB2Gray--Reduce--Sobel work reduced principal storage approximately from
\[
M_{\rm separate}\approx 11N
\qquad\text{to}\qquad
M_{\rm composed}\approx 7N.
\]
The largest successful square side increased from \(87{,}552\) to \(109{,}824\), corresponding
to a \(57.3\%\) increase in feasible pixel count.  At \(2048^2\), however, the composed kernel was
about \(24\%\) slower because grayscale values were recomputed.  This is exactly the kind of
tradeoff the theory must expose: storage reduction is a provable semantic/resource consequence;
execution time follows from the selected realization and machine properties.

\section{The Existing Attention Family Is Part of the Same Theory}

The attention papers are not external kernels to be wrapped by a new Transformer theory.  They
are prior MoA derivations to be composed.

The established progression is
\[
\boxed{\begin{aligned}
\text{forward attention}&\rightarrow\text{backward/fused training}\\
&\rightarrow\text{decode/KV cache/GQA/MQA}\rightarrow\text{complete pre-norm block}.
\end{aligned}}
\]

Forward attention uses systematic \(\psi\) reduction to access \(Q,K,V\) directly, eliminating an
explicit \(K^T\) buffer and unnecessary score/softmax intermediates.  Backward attention applies
the same discipline to the softmax VJP and gradients.  Decode specializes the DNF to a current
query row, expresses KV-cache accumulation with canonical catenation \(\#\), and expresses
GQA/MQA through \(\psi\) selection.  The normalization/FFN work extends the same method to
RMSNorm/LayerNorm and gated MLP/SwiGLU.

Thus the Transformer work below must reuse those derivations rather than recreate conventional
tensor formulas and translate them afterward.

\section{Complete Transformer Block: Architecture-Independent Formulation}

Let the input to block \(l\) be
\[
\boxed{
\rho X^{(l-1)}=\langle B\;N\;D\rangle
}
\]
and the output satisfy
\[
\boxed{
\rho X^{(l)}=\langle B\;N\;D\rangle.
}
\]

For the pre-norm block, introduce names only for expressions that will subsequently be substituted:
\[
Z^{(l)}=\Norm_1^{(l)}(X^{(l-1)}),
\]
\[
A^{(l)}=\Attn^{(l)}(Z^{(l)}),
\]
\[
R^{(l)}=X^{(l-1)}+A^{(l)},
\]
\[
W^{(l)}=\Norm_2^{(l)}(R^{(l)}),
\]
\[
F^{(l)}=\FFN^{(l)}(W^{(l)}),
\]
and
\[
\boxed{X^{(l)}=R^{(l)}+F^{(l)}.}
\]

These names do not assert materialized arrays.  They are substitution points on the path to DNF.

\subsection{Block boundary in \(\psi\) notation}

For every
\[
\langle b\#n\#d\rangle\in\iota(\rho X^{(l)}),
\]
\[
\boxed{
\begin{aligned}
\langle b\#n\#d\rangle\psi X^{(l)}
={}&
\langle b\#n\#d\rangle\psi X^{(l-1)}
\\
&+
\langle b\#n\#d\rangle\psi A^{(l)}
\\
&+
\langle b\#n\#d\rangle\psi F^{(l)}.
\end{aligned}}
\]

This is the component from which the complete block DNF is derived.

\section{Normalization Defined by Shape}

To avoid collision with the fundamental translator \(\gamma\), this working document denotes the
learned RMSNorm scale by \(\vec g_1^{(l)}\) and \(\vec g_2^{(l)}\).  The symbol \(\gamma\) is
reserved for the DNF-to-ONF address/layout family.

For fixed \(\langle b\#n\rangle\), define the first normalization statistic
\[
\boxed{
r^{(l)}_{1,b,n}
=
\sqrt{
\frac{1}{D}
\sum_{e\in\iota(D)}
\left(
\langle b\#n\#e\rangle\psi X^{(l-1)}
\right)^2+\epsilon
}.
}
\]
Then
\[
\boxed{
\langle b\#n\#d\rangle\psi Z^{(l)}
=
g^{(l)}_{1,d}
\frac{
\langle b\#n\#d\rangle\psi X^{(l-1)}
}{
r^{(l)}_{1,b,n}
}.
}
\]

The residual component is
\[
\boxed{
\langle b\#n\#d\rangle\psi R^{(l)}
=
\langle b\#n\#d\rangle\psi X^{(l-1)}
+
\langle b\#n\#d\rangle\psi A^{(l)}.
}
\]

The second normalization statistic is therefore
\[
\boxed{
r^{(l)}_{2,b,n}
=
\sqrt{
\frac{1}{D}
\sum_{e\in\iota(D)}
\left(
\langle b\#n\#e\rangle\psi R^{(l)}
\right)^2+\epsilon
}.
}
\]
and
\[
\boxed{
\langle b\#n\#d\rangle\psi W^{(l)}
=
g^{(l)}_{2,d}
\frac{
\langle b\#n\#d\rangle\psi R^{(l)}
}{
r^{(l)}_{2,b,n}
}.
}
\]

All bounds arise from shapes.  No machine dimension has been introduced.

\section{FFN Defined by Shape and \(\psi\)}

Let
\[
\rho W_{\rm gate}^{(l)}
=
\rho W_{\rm up}^{(l)}
=
\langle D\;D_{\rm ff}\rangle,
\qquad
\rho W_{\rm down}^{(l)}
=
\langle D_{\rm ff}\;D\rangle.
\]

For
\[
f\in\iota(D_{\rm ff}),
\]
define
\[
u^{(l)}_{b,n,f}
=
\sum_{d'\in\iota(D)}
(\langle b\#n\#d'\rangle\psi W^{(l)})
\times
(\langle d'\#f\rangle\psi W_{\rm gate}^{(l)}),
\]
\[
v^{(l)}_{b,n,f}
=
\sum_{d'\in\iota(D)}
(\langle b\#n\#d'\rangle\psi W^{(l)})
\times
(\langle d'\#f\rangle\psi W_{\rm up}^{(l)}),
\]
and
\[
h^{(l)}_{b,n,f}
=
\SiLU(u^{(l)}_{b,n,f})
\times
v^{(l)}_{b,n,f}.
\]

Then
\[
\boxed{
\langle b\#n\#d\rangle\psi F^{(l)}
=
\sum_{f\in\iota(D_{\rm ff})}
h^{(l)}_{b,n,f}
\times
(\langle f\#d\rangle\psi W_{\rm down}^{(l)}).
}
\]

Substitution removes \(U,V,H\) as required semantic intermediates:
\[
\boxed{
\begin{aligned}
\langle b\#n\#d\rangle\psi F^{(l)}
={}&
\sum_{f\in\iota(D_{\rm ff})}
\Bigg[
\SiLU\!\left(
\sum_{d'\in\iota(D)}
(\langle b\#n\#d'\rangle\psi W^{(l)})
(\langle d'\#f\rangle\psi W_{\rm gate}^{(l)})
\right)
\\
&\quad\times
\left(
\sum_{d'\in\iota(D)}
(\langle b\#n\#d'\rangle\psi W^{(l)})
(\langle d'\#f\rangle\psi W_{\rm up}^{(l)})
\right)
\\
&\quad\times
(\langle f\#d\rangle\psi W_{\rm down}^{(l)})
\Bigg].
\end{aligned}}
\]

The remaining \(W^{(l)}\) occurrences are then replaced by the second-normalization component
definition, and \(R^{(l)}\) is replaced by its residual component definition.

\section{Attention Is Substituted from the Established MoA DNF}

The next step is not to invent another attention formula.  Every occurrence of \(Z^{(l)}\) in the
established attention-family DNF is replaced by the first-normalization \(\psi\) expression above.
The attention result \(A^{(l)}\) is then substituted into both the residual and the FFN path.

The derivation therefore proceeds
\[
\boxed{\begin{aligned}
\text{Transformer formulation}
&\rightarrow \text{established MoA attention DNF}\rightarrow \text{residual/norm substitution}\\
&\rightarrow \text{FFN substitution}\rightarrow \psi\text{ reduction}\rightarrow \text{candidate block DNF}.
\end{aligned}}
\]

The candidate DNF is frozen only after component, shape, and numerical verification.

\section{Transformer Invariants and Verification Gate}

The first block invariant is
\[
\boxed{
\rho X^{(l)}
=
\rho X^{(l-1)}
=
\langle B\;N\;D\rangle.
}
\]

Thus blocks compose:
\[
X^{(0)}
\xrightarrow{B^{(1)}}
X^{(1)}
\xrightarrow{B^{(2)}}
\cdots
\xrightarrow{B^{(L)}}
X^{(L)},
\]
with
\[
\boxed{
\rho X^{(0)}=\rho X^{(1)}=\cdots=\rho X^{(L)}
=
\langle B\;N\;D\rangle.
}
\]

At every transformation boundary the verification rule is
\[
\boxed{
\text{same semantic function}
+
\text{same result shape}
+
\text{same components}.
}
\]

After dimension lifting, add the execution-space invariant
\[
\boxed{\tau(\rho P^L)=\tau(\rho P).}
\]

Therefore the complete gate is
\[
\boxed{
\text{derive}
\rightarrow
\text{verify semantics}
\rightarrow
\text{verify shape}
\rightarrow
\text{verify components}
\rightarrow
\text{FREEZE}
\rightarrow
\text{continue}.
}
\]

Only after the complete Transformer-block DNF passes this gate do we introduce a chosen
\(\gamma\), sequential ONF, dimension lifting, and a target machine array.

\section{The Unified Scientific Test}

The CPU, GPU, multi-GPU, distributed, CPU--GPU, attention, decode, and Transformer experiments
must validate one another in the following precise sense: each is an independent instantiation of
the same theory and therefore tests a different implication of the same invariant definitions.

\[
\boxed{
\begin{array}{c}
\text{pointwise}\\
\text{neighborhood}\\
\text{reduction}\\
\text{composition}\\
\text{attention}\\
\text{Transformer}
\end{array}}
\quad
\xrightarrow{\rho,\iota,\psi}
\quad
\boxed{\text{DNF}}
\quad
\xrightarrow{\gamma,\;\text{reshape},\;\rav}
\quad
\boxed{
\begin{array}{c}
\text{CPU}\\
\text{GPU}\\
\text{multi-GPU}\\
\text{distributed}\\
\text{CPU--GPU}
\end{array}}
\]

A contradiction between independently derived experiments is therefore scientifically useful:
it identifies either an incorrect semantic derivation, an incorrect machine description, an
incorrect mapping, or an incomplete machine-cost property.  It must not be repaired by changing
the theory independently for each backend.

Conversely, agreement across independently frozen experiments strengthens the common theory
because the same definitions survive changes in algorithm class, dimensionality, machine shape,
and realization.

\section{Next Mathematical Step}

The immediate next step is to complete the substitution
\[
\boxed{\langle b\#n\#d\rangle\psi X^{(l)}}
\]
using the established MoA attention-family DNF until only true block inputs, weights, scalar
operations, reductions, \(\rho\), \(\iota\), and \(\psi\) remain.

At that point:
\[
\boxed{\text{complete Transformer-block DNF}}
\]
is verified and frozen.

Only afterward:
\[
\boxed{\begin{aligned}
\DNF &\xrightarrow{\gamma_{\rm chosen}} \text{Sequential ONF}
\xrightarrow{\text{verify}} \text{Dimension Lift}\\
&\xrightarrow{\text{map}} \text{Machine Array}
\xrightarrow{\text{predict/freeze}} \text{execute/validate}.
\end{aligned}}
\]

The same procedure then scales from one block to multiple blocks and eventually to complete
models without changing the underlying mathematics.

\section*{Working Principle}

\[
\boxed{\textbf{Define by shape. Generate indices with \(\iota\). Select with \(\psi\).}}
\]
\[
\boxed{\textbf{Reduce to one invariant DNF.}}
\]
\[
\boxed{\textbf{Translate with an explicit \(\gamma\), reshape/lift, and \(\rav\).}}
\]
\[
\boxed{\textbf{Treat the machine as an array. Predict, freeze, execute, validate.}}
\]
\[
\boxed{\textbf{One theory. Independent experiments. Mutual validation. MoA all the way.}}
\]

\begin{thebibliography}{99}

\bibitem{mullin1988}
L. M. Mullin,
\emph{A Mathematics of Arrays},
Ph.D. dissertation, Syracuse University, 1988.

\bibitem{mullinjenkins1996}
L. M. Mullin and M. A. Jenkins,
``Effective Data Parallel Computation Using the Psi Calculus,''
\emph{Concurrency: Practice and Experience}, 8(7), 499--515, 1996.

\bibitem{mullin2005}
L. Mullin,
``A Uniform Way of Reasoning about Array-Based Computation in Radar:
Algebraically Connecting the Hardware/Software Boundary,''
\emph{Digital Signal Processing}, 15(5), 466--520, 2005.

\bibitem{attentionforward}
L. M. Mullin and G. Hains,
``Attention at the Theoretical Minimum: A Mathematics of Arrays Framework
for Memory-Optimal Transformer Kernels,'' 2026.

\bibitem{attentionbackward}
L. M. Mullin and G. Hains,
``Attention at the Theoretical Minimum: Backward Pass MoA DNF Derivation,
PyTorch Implementation, and Numerical Verification,'' 2026.

\bibitem{decode}
L. M. Mullin and G. Hains,
``MoA-Structured Decode Attention: DNF Derivation, KV-Cache Accumulation,
GQA/MQA, and OpenACC Kernel,'' 2026.

\bibitem{transformerblock}
L. M. Mullin and G. Hains,
``The Transformer Block at the Theoretical Minimum: Memory-Optimal
Normalization and Feed-Forward Layers via the Mathematics of Arrays,'' 2026.

\bibitem{compiler}
L. M. Mullin,
\emph{MoA Compiler Research Repository},
\url{https://github.com/womenflyplanes/moa-compiler}, 2026.

\end{thebibliography}

\section*{Acknowledgment}
The author acknowledges the computational resources and research-computing support used across
the CPU, GPU, multi-GPU, distributed, and heterogeneous CPU--GPU experiments, including
University at Albany and ACCESS-supported resources where applicable.  These independent
machine realizations provide the experimental basis for testing one common MoA theory.


%% === DELTA REAL HARDWARE PROOF - Added 2026-10-03 ===
%% NCSA Delta bihf-delta-gpu 2993hrs - Nodes: dt-login04 x86_64, gpua027 A100, gpub017 A40
%% All 12 files same F M T proves Same DNF -> Many ONFs -> Predictable
\begin{verbatim}
Fortran gfortran col-major gamma_col=7:
 paper1_exact F=      786432  M=         768  T ms=   6.21772837E-03
C row-major gamma_row=6:
 paper1_exact F=786432 M=768 T=0.006218 ms bytes=3.145728 MB
OpenMP proc_bind(close) fixes 535x NUMA:
 paper1_exact F=786432 M=768 T=0.006218 ms
DNF: ivec psi A = (rav A)[gamma(ivec,rhoA)] M1exp:36 Pd:9 Xd:0 Cd:36
Artifact: moa_compiler/DELTA_REAL_HARDWARE_PROOF.txt moa_A100_22646116.log moa_A40_22646132.log
\end{verbatim}
%% === END DELTA PROOF ===

\end{document}
